% !TeX root=../main.tex
\chapter{پیاده‌سازی}

% TODO ۴-۱ مقدمه فصل — طبق بند ۱-۷ قالب: این فصل محیط و ابزارهای اجرا، معیار ارزیابی
% نتایج، ارائه و تحلیل نتایج، و ویژگی‌ها و دستاوردهای روش پیشنهادی را می‌گوید. یک جمله
% هم درباره‌ی مرز فصل ۳ و ۴ (فصل ۳ روش و پیاده‌سازی، فصل ۴ اجرا و ارزیابی).
\section{مقدمه}

\section{محیط و ابزارهای اجرا}

% TODO ۴-۲-۱ مشخصات ماشین محلی (پردازنده، هسته، حافظه، سیستم‌عامل)؛ مسیر پردازنده‌ی
% گرافیکی ابری؛ نسخه‌ی پایتون؛ محیط مجازی و فایل قفل وابستگی (چرا pip freeze و نه بازه‌ی
% نسخه). ⚠️ مشخصات سخت‌افزاری باید از سیستم واقعی خوانده شود، نه حدس.
% منبع: ردیف ۱ doc/decision_log.md؛ requirements.lock؛ Makefile
\subsection{محیط سخت‌افزاری و نرم‌افزاری}

% TODO ۴-۲-۲ جدول (۴-۱): کتابخانه‌های اصلی، نسخه، و نقش هرکدام در پروژه (فقط موارد
% مؤثر بر نتیجه، نه کل requirements.lock).
% منبع: requirements.lock
\subsection{کتابخانه‌های اصلی}

\subsection{زیرساخت ثبت آزمایش و اجرای موازی}

ثبت آزمایش با \lr{MLflow} روی یک \lr{backend} فایلی محلی انجام شد. تا پایان فاز ۷،
۲٬۲۱۲ اجرا ثبت شده است؛ مقایسه‌ی دستی این حجم غیرقابل‌دفاع می‌شد، پس هر اجرا موظف شد
اطلاعات لازم برای بازیابی و مقایسه را خودش حمل کند. نخستین لایه‌ی این تعهد در
\textbf{نام} اجراست: قالبی ثابت که خانواده، مدل، سطح داده، متغیر هدف، مجموعه‌ی ویژگی،
نقطه‌ی $\tau$، مرحله، سید و مُهر زمانی را به‌ترتیب کنار هم می‌گذارد، چنان‌که کل مختصات
یک آزمایش از خودِ نام خوانده شود. لایه‌ی دوم فیلدهای اجباری است؛ در کنار برچسب‌های
مرحله و منبع محاسبه و خانواده، دو برچسب اختصاصی ثبت می‌شوند: \lr{\texttt{model\_type}}
(کلاس یا کتابخانه‌ی واقعی الگوریتم) و \lr{\texttt{code\_ref}} (مسیر فایل، شماره‌ی خط و
نام تابع آموزش‌دهنده، که خودکار از خود تابع استخراج می‌شود). کنار برچسب \lr{commit}
گیت، این یعنی بازتولید یک اجرا به \lr{checkout} کردن همان \lr{commit} و بازکردن دقیقاً
همان آدرس خلاصه می‌شود. پارامترها شامل هش \lr{snapshot} داده، هش \lr{fold}ها، شناسه‌ی
مجموعه‌ی ویژگی، سطح داده و سیدند، و معیارها علاوه بر زمان برازش شامل معیارهای عملیاتی
کامل‌اند (تفصیل معیارها در بند ۴-۳).

مهم‌ترین قید این زیرساخت \textbf{دروازه‌ی انصاف} است. فایل \lr{fold}ها یک بار ساخته شد،
هشش در دفترچه‌ی داده ثبت و فایل منجمد اعلام شد (رویه‌ی ساختش در بند ۳-۴-۳ آمده). هر
اجرایی با هش \lr{fold} نامنطبق از جدول مقایسه حذف می‌شود؛ بدون این قید هیچ راهی برای
اثبات اینکه یک اجرای ابری واقعاً روی همان برش‌های زمانی انجام شده وجود نداشت.

اجرای محلی اولویت اول بود، اما چند خانواده روی پردازنده زمان غیرقابل‌قبول داشتند:
شبکه‌ی عصبی روی سطح رزرو فردی با بیش از دو میلیون رکورد، فرآیند گاوسی روی ماتریسی به
ابعاد ۷٬۵۷۹ در ۷٬۵۷۹، و نمونه‌گیری بیزی سنگین. برای این‌ها مسیر پردازنده‌ی گرافیکی
رایگان طراحی شد، با قیدهایی که همان قاعده‌ی «کد در \lr{\texttt{src/}} است» را حفظ
می‌کنند: کد اصلی در نوت‌بوک نوشته نمی‌شود و یک هدف \lr{\texttt{Makefile}} کل
\lr{\texttt{src/}} را با داده‌ی لازم بسته‌بندی می‌کند؛ سلول سوم هر نوت‌بوک روی هش
\lr{snapshot} و هش \lr{fold}ها \lr{assert} می‌زند و در صورت عدم تطابق متوقف می‌شود، چون
نوت‌بوک ابری داده را از فضای ذخیره‌سازی شخصی می‌خواند نه از \lr{DVC}؛ وضعیت بهینه‌سازی
هر ده \lr{trial} ذخیره می‌شود تا اجرا پس از قطع نشست ادامه‌پذیر باشد؛ و نوت‌بوک اجراشده
\textbf{حتی در صورت شکست آزمایش} به مخزن برمی‌گردد، چون شکست هم داده است. هر چهار
نوت‌بوک برنامه‌ریزی‌شده اجرا شدند و برگشتند.

چند محور از این پروژه به‌طور طبیعی مستقل‌پذیرند: هر خانواده‌ی مدل، هر سری زمانی درون
خانواده‌ی سری‌ها، هر نقطه‌ی $\tau$ و هر سید. قاعده‌ی حاکم بر اجرای موازی یک جمله است:
\textbf{هیچ دو اجرای هم‌زمان یک فایل مشترک را نمی‌نویسند}. هر اجرا نام یکتای خودش را
می‌سازد و خروجی‌اش را زیر همان نام می‌نویسد، و مسیرهای خروجی اجراهای ابری از مسیرهای
محلی جدا نگه داشته شدند. اثر عملی این قاعده قابل اندازه‌گیری بود: مرحله‌ی غربالگری
نخستین خانواده سیصد کار مستقل بود که روی هشت پردازه‌ی موازی در حدود پنجاه‌ونه دقیقه
تمام شد، و چهار نوت‌بوک پردازنده‌ی گرافیکی هم‌زمان روی چهار حساب کاربری مستقل اجرا
شدند. استثنا هم صریح اعلام شد: خانواده‌ی ترکیب مدل‌ها و لایه‌ی کالیبراسیون
\textbf{وابسته}‌اند و فقط پس از تعیین قهرمانان اجرا می‌شوند.

\section{معیار ارزیابی}

% TODO ۴-۳-۱ pinball@tau (فضای نرخ و تعداد پرس)، coverage@tau، shortage\_rate،
% waste\_reduction\_pct، RMSE/MAE؛ چرا معیار اصلی pinball است (ارجاع به بند ۲-۲-۲).
% منبع: doc/model-evaluation-metrics.md
\subsection{مجموعه‌ی معیارها}

% TODO ۴-۳-۲ پنجره‌ی Test قفل‌شده؛ اعتبارسنجی rolling-origin (ارجاع به بند ۲-۴-۱)؛
% ارجاع به بند ۳-۴-۳ برای رویه‌ی ساخت فولدها، اینجا فقط منطق پروتکل.
% منبع: doc/progress/06-*.md
\subsection{پروتکل اعتبارسنجی}

% TODO ۴-۳-۳ خط پایه‌های B1-B6 و نقش B3 به‌عنوان رقیب واقعی.
% منبع: reports/baselines.md
\subsection{خط پایه‌ها}

% TODO ۴-۳-۴ Diebold-Mariano؛ بوت‌استرپ بلوکی دوبعدی و n\_eff؛ تصحیح چندگانه (طبق
% بند ۲-۵).
\subsection{آزمون معناداری}

\section{نتایج به‌دست‌آمده}

% TODO ۴-۴-۱ جدول (۴-۲): مقایسه با CI بوت‌استرپ بلوکی دوبعدی (فقط ۸-۱۰ ردیف شاخص،
% جدول کامل به پیوست)؛ نتایج منفی و پرونده‌های بسته‌شده گزارش شوند (معماری دومرحله‌ای،
% خوشه‌بندی شهر).
% منبع: ردیف ۴۰، ۴۴-۴۸ doc/decision_log.md
\subsection{نتایج اعتبارسنجی و انتخاب نامزدها}

% TODO ۴-۴-۲ ⚠️ مهم‌ترین صداقت گزارش: قهرمان اعتبارسنجی و قهرمان Test یکی نیستند، با
% CI و توضیح علت، نه پنهان‌کاری. تأکید بر «Test فقط یک بار باز شد».
% منبع: ردیف ۴۹ doc/decision_log.md؛ reports/phase8/
\subsection{ارزیابی نهایی روی مجموعه‌ی قفل‌شده}

% TODO ۴-۴-۳ پوشش حاشیه‌ای در برابر شرطی؛ شکاف نرخ کمبود مشاهده‌شده تا tau اسمی.
% منبع: reports/phase8/
\subsection{کالیبراسیون و پوشش}

\section{تحلیل نتایج}

% TODO ۴-۵-۱ تحلیل باقیمانده؛ تجزیه‌ی خطا؛ الگوی بدترین موارد؛ پایداری/حساسیت؛
% Overfit/Underfit.
% منبع: reports/phase8/
\subsection{تحلیل خطا و بدترین موارد}

% TODO ۴-۵-۲ اهمیت ویژگی و SHAP؛ اثرات نهایی (PDP/ALE)؛ مدل جانشین قابل‌فهم.
% منبع: reports/phase9/
\subsection{تفسیرپذیری}

% TODO ۴-۵-۳ تبیین تناقض ظاهری log\_res (کوانتایل در برابر میانگین، طبق بند ۲-۳-۲).
% منبع: ردیف ۵۱ doc/decision_log.md
\subsection{اعتبارسنجی دامنه‌ای}

% TODO ۴-۵-۴ نتیجه‌ی L1/L2 در برابر L5 با CI.
% منبع: ردیف ۵۰ doc/decision_log.md
\subsection{مقایسه‌ی سطح تجمیعی در برابر سطح فرد}

% TODO ۴-۵-۵ منحنی trade-off کمبود/هدررفت؛ سناریوهای سیاستی؛ صرفه‌جویی سالانه با CI
% (نه عدد نقطه‌ای)؛ پروتکل پایلوت.
% منبع: reports/phase10/
\subsection{تحلیل کسب‌وکاری}

\section{ویژگی‌ها و دستاوردهای روش پیشنهادی}

% TODO ۴-۶-۱ مزیت‌های عملی روش: کارکرد در سطح عملیاتی (d,m,r,f)؛ کنترل‌پذیری صریح
% نقطه‌ی tau به‌عنوان اهرم سیاستی؛ صرفه‌جویی برآوردی با CI. ⚠️ فقط ادعاهایی که در بند
% ۴-۴ شاهد عددی دارند.
% منبع: reports/phase10/؛ ردیف ۳۶ doc/decision_log.md
\subsection{مزیت‌های عملی}

\subsection{بازتولیدپذیری}

آنچه در این پروژه قفل شده چهار چیز است: نسخه‌ی داده (هش \lr{SHA-256} هر فایل خام و
پردازش‌شده)، برش‌های اعتبارسنجی (فایل منجمد \lr{fold}ها و هشش، که هر اجرای نامنطبق را
از مقایسه حذف می‌کند)، سید تصادفی (منبع واحد در ماژول پیکربندی)، و مرجع کد هر نتیجه
(\lr{commit} گیت به‌همراه آدرس دقیق فایل و خط تابع آموزش‌دهنده، هر دو در برچسب‌های همان
اجرای \lr{MLflow}).

آنچه قفل \textbf{نشده} هم باید صریح گفته شود. مخزن \lr{DVC} فقط \lr{cache} محلی دارد و
مقصد راه دور برایش تنظیم نشده، یعنی بازتولید کامل روی ماشین دیگر نیازمند انتقال دستی
داده‌ی خام است؛ مخزن گیت هم مقصد راه دور ندارد و رساندن کد به محیط ابری از مسیر بسته‌ی
فشرده انجام شد. افزون بر این، بازتولید سرتاسری فعلاً زنجیره‌ای از دستورهای ماژولی است و
یک هدف واحد که کل مسیر داده‌ی خام تا جدول نتایج را با یک فرمان اجرا کند هنوز ساخته
نشده است.

% TODO: اگر بسته‌کاری ۱۱-۴ (README + هدف واحد `make all`) اجرا شد، پاراگراف بالا
% به‌روزرسانی شود. منبع: doc/progress/11-*.md بند ۱۱.۴

% TODO ۴-۶-۳ ⚠️ حداکثر نیم صفحه (تصمیم کاربر: فقط یک اشاره‌ی کوتاه). محتوا: تقسیم کار
% «انسان صاحب تصمیم و پذیرش، عامل مجری و مستندساز»؛ قرارداد قابل‌بارگذاری AGENTS.md؛
% چهار دفتر دائمی و ضمانت مکانیکی src/report_coverage.py با کد خروجی ناصفر؛ و نتیجه‌ی
% عملی: هرجا قرارداد به قید کد تبدیل شد رعایت شد، هرجا فقط سند ماند نقض شد. جدول چهار
% اشتباه فاز ۷ و باگ سیم‌کشی فیچرست به پیوست می‌روند، نه به بدنه.
% منبع: AGENTS.md؛ doc/phase7-execution-standard.md؛ ردیف ۳۷، ۴۷ doc/decision_log.md
\subsection{اجرای پروژه با کمک عامل‌های هوش مصنوعی}

% TODO ۴-۷ خلاصه و جمع‌بندی فصل
\section{خلاصه و جمع‌بندی}
